% Hurricane Milton case-study summary table. Regenerated 2026-06-11 from
%   /iopsstor/scratch/cscs/sadamov/milton_case_study/milton_per_baseline_summary.csv
% via /iopsstor/scratch/cscs/sadamov/milton_case_study/analyze_milton_stats.py.
% Post-hoc rows updated to Phase 5 (IFS-ENS perturbed-IC + weight) configuration
% per the post-Phase-5 narrative decision (see [[phase5-milton-redo-decision]]).
% 8 baselines x 14 inits x 10 members = 140 (init, member) pairs per row.
% Detection rate = fraction of those 140 that produced any track within the
% Milton lifecycle window (2024-10-04 .. 2024-10-11). Position errors are
% mean great-circle distance against IBTrACS v4 NA-basin best track and
% against the 0.25 deg ERA5 control track, averaged over all matched
% (init, member, lead) rows. Intensity bias is mean MSL minus IBTrACS MSL
% at matched validating times, in hPa; the +37 to +44 hPa cluster reflects
% the 0.25 deg ERA5 resolution-limited intensity ceiling. Weight-only
% post-hoc numbers preserved in the v0 commit (cb6a156 paper repo) for the
% data release.
%
% IFS-ENS row provenance. On 2026-09-01 this row was changed from
% 71 / 317 / 173 / +36.6 / 252 to 14 / 348 / 169 / +40.2 / 287 after an audit
% could not reproduce it. That audit was wrong and the change was reverted on
% 2026-09-02: aggregate_tracks.py carried a stale double member filter. The
% tracker (track_one_init.py) already selects the stratified IFS-ENS 10-member
% subset and relabels it 0..9 like every other baseline, and the aggregator then
% re-applied the original {0,5,...,45} stratification to those relabelled ids,
% which matches only {0,5}. It therefore scored 2 of 10 members against a
% denominator of 140, giving 19/140 = 14%; counted over all ten members the same
% archived tracks give 99/140 = 71%. The filter was gated on baseline=="ifs_ens",
% so no ML row was ever affected, which is why only this row looked anomalous.
% The audit's apparent corroboration (job 2473058, ~12.6%) predates the
% --maxgap 6h -> 18h change and is a separate, since-fixed artefact.
% Filter removed 2026-09-02; the pipeline now reproduces this row to four
% significant figures (71% / 316.5 / 172.7 / +36.55 / 251.6).
%
% Known and unrepaired here: ~20% of (variable, init, lead) chunks in the
% WeatherBench-2 IFS-ENS archive hold a globally NaN field, independently per
% variable, leaving only 41% of leads with all six tracked fields present (vs
% ~100% for the ML baselines). Repairing that by linear interpolation along
% lead_time raises IFS-ENS to 80% detection / 326 km / +33.6 hPa and leaves the
% ranking unchanged; the unrepaired 71% is reported here as the conservative
% choice. See .claude/logbook/2026-09-02_milton_ifsens_detection_double_filter_bug.md.

\begin{table}[t]
  \centering
  \small
  \setlength{\tabcolsep}{6pt}
  \caption{Hurricane Milton (October 2024) case-study summary. Each row
           aggregates 14 init times $\times$ 10 ensemble members. Detection
           rate is the fraction of those 140 (init, member) pairs that
           produced any track point within the Milton lifecycle window.
           Position errors are mean great-circle distance, averaged over
           all matched (init, member, lead) rows. The IBTrACS-axis MSL
           bias is dominated by the 0.25$^\circ$ ERA5 resolution-limited
           intensity ceiling and is therefore not a baseline-quality
           discriminator (see text). Sorted by detection rate descending
           within each role group. All baselines, including IFS-ENS,
           evaluate on the same stratified 10-member ensemble (140 cells
           per baseline); the tracker bridges the WeatherBench-2 IFS-ENS
           archive gaps (sec.~\ref{sec:methods-metrics}) with an 18\,h
           maximum-gap setting, chosen to match the longest contiguous
           archive gap; a tighter limit would break IFS-ENS tracks at
           those gaps (an archive artefact, not forecast skill), while
           the ML tracker inputs carry no such gaps. IFS-ENS's rate
           (71\%) is a lower bound: the archive stores roughly a fifth of
           each field's lead times as missing, independently per variable,
           so only 41\% of its leads carry all six tracked fields against
           $\sim$100\% for the ML baselines; reconstructing those leads
           by interpolation raises it to 80\% and leaves the ordering
           unchanged. \textbf{Bold} marks the best value in each
           scored column (highest detection rate, lowest position error); the
           MSL-bias column is not bolded, being resolution-limited rather than a
           baseline-quality discriminator (see text).}
  \label{tab:milton-summary}
  \begin{tabular}{@{}l l rrrrr@{}}
    \toprule
    \textbf{Role} & \textbf{Model}
      & \textbf{Detect.}
      & \textbf{Mean pos.}
      & \textbf{Median pos.}
      & \textbf{Mean MSL}
      & \textbf{Mean pos.} \\
     &
      & \textbf{rate (\%)}
      & \textbf{err. (km)}
      & \textbf{err. (km)}
      & \textbf{bias (hPa)}
      & \textbf{err. vs ERA5 (km)} \\
    \midrule
    \rowcolor{phaseBg}
    \multicolumn{7}{@{}l}{\textbf{Trained probabilistic}} \\
    & Atlas         & \textbf{89} & 281 & 159 & $+42.4$ & 283 \\
    & FCN3          & 87 & 438 & 225 & $+44.2$ & 435 \\
    & AIFS-ENS      & 86 & 300 & 163 & $+39.2$ & 300 \\
    \hline
    \rowcolor{phaseBg}
    \multicolumn{7}{@{}l}{\textbf{IC-augmented SPW (perturbed-IC + weight, this work)}} \\
    & sfno\_modes10\_ic    & 76 & 306 & 161 & $+41.7$ & 254 \\
    & graphcast\_all\_ic   & 73 & 258 & 157 & $+40.0$ & 250 \\
    & aifs\_perturbed\_ic  & 67 & 262 & 133 & $+40.3$ & 245 \\
    & aurora\_encoder\_ic  & 66 & \textbf{218} & \textbf{124} & $+38.6$ & \textbf{191} \\
    \hline
    \rowcolor{phaseBg}
    \multicolumn{7}{@{}l}{\textbf{Classical reference}} \\
    & IFS-ENS          & 71 & 317 & 173 & $+36.6$ & 252 \\
    \bottomrule
  \end{tabular}
\end{table}
